\documentclass[conference]{IEEEtran}
\IEEEoverridecommandlockouts
% The preceding line is only needed to identify funding in the first footnote. If that is unneeded, please comment it out.
\usepackage{cite}
\usepackage{amsmath,amssymb,amsfonts}
\usepackage{algorithmic}
\usepackage{graphicx}
\usepackage{textcomp}
\usepackage{xcolor}
\usepackage{booktabs}
\usepackage{multirow}
\usepackage{array}
\def\BibTeX{{\rm B\kern-.05em{\sc i\kern-.025em b}\kern-.08em
    T\kern-.1667em\lower.7ex\hbox{E}\kern-.125emX}}
\begin{document}

\title{Machine Learning-Based Flood Prediction Using Sequential Weather Patterns: An SVM–LSTM Comparative Study\thanks{No external funding was used for this work.}}

\author{

\IEEEauthorblockN{1\textsuperscript{st} V. Prajesh}
\IEEEauthorblockA{
\textit{Dept. of Computer Science Engineering (AI \& ML)} \\
\textit{Bangalore, India} \\
prajeshv.03@gmail.com
}

\and

\IEEEauthorblockN{2\textsuperscript{nd} Ankith D Prabhu}
\IEEEauthorblockA{
\textit{Dept. of Computer Science Engineering (AI \& ML)} \\
\textit{Bangalore, India} \\
ankithprabhu46@gmail.com
}

\and

\IEEEauthorblockN{3\textsuperscript{rd} Yisol Choi}
\IEEEauthorblockA{
\textit{Dept. of Computer Science Engineering (AI \& ML)} \\
\textit{Bangalore, India} \\
yisolchoi04@gmail.com
}

\and

\IEEEauthorblockN{4\textsuperscript{th} Shreya H Rao}
\IEEEauthorblockA{
\textit{Dept. of Computer Science Engineering (AI \& ML)} \\
\textit{Bangalore, India} \\
raoshreya018@gmail.com
}

}


\maketitle

\begin{abstract}
Abstract—Short-term flood forecasting is very essential in disaster
risk reduction in flood-prone regions like India, where high intensity
monsoon rainfall, complex river basins, and rapid urbanization
The frequency and incidence of flood events will simultaneously increase. Recent
Through advances in machine learning, there are data-driven models that
to complement traditional hydrological modeling, especially for
Operational forecasting at the daily or sub-daily time scale. This
Paper presents comparative analysis of the Support Vector Machine.
(SVM) and Long Short-Term Memory (LSTM) neural network
models for short-term flood occurrence prediction using hydro-
meteorological time series derived from the publically available
Flood Risk in India dataset.
Using a subset of 10,000 daily records, each containing
Rainfall, river discharge, temperature, humidity and a binary
flood label, from which supervised 7-day input sequences are
constructed to predict the occurrence of flood on the next day. Data are
split chronologically into 80% training and 20% testing, and all
Continuous features are scaled with a Min–Max scaler fit on the
training set to avoid data leakage. The SVM with a radial basis
function kernel operates on flattened 7-day windows, while
the LSTM processes the full sequence using a 64-unit recurrent
layer followed by dense layers with dropout regularization, and
Both models include class weights derived from training.
distribution.
Model performance is measured based on accuracy, precision,
recall, F1-score, confusion matrices, and area under the receiver
operating characteristic curve (ROC–AUC) on the held-out
test set of 1,999 samples. The SVM achieves an accuracy of
0.4987, F1-score of 0.4898, and ROC–AUC of 0.5069, indicating
achieved only marginally better-than-random discrimination. The LSTM
achieves a slightly better accuracy of 0.5038, whereas it is substantially higher
flood-class F1-score of 0.6653, and extremely high recall of
0.9860, but with ROC–AUC close to random (0.5013) and a
large number of false alarms. These findings indicate that, under
the current feature set and configuration, the LSTM outperforms
prioritises the detection of flood-events at the expense of specificity while
Both models need richer inputs, refined architectures, and
threshold calibration for reaching operationally robust performance.
\end{abstract}

\begin{IEEEkeywords}
flood prediction, hydro-meteorological time series, Support Vector Machine (SVM), Long Short-Term Memory (LSTM), machine learning, India.
\end{IEEEkeywords}

\section{Introduction}
Floods are one of the most common and damaging natural
hazards worldwide, and India is particularly vulnerable due to
to its monsoon-dominated climate, dense river network, and
rapidly urbanized floodplains. Recurring flood events cause
Substantial loss of life, damage to infrastructure, agricultural
losses, and extreme interruptions to livelihoods and regional
economies. In such a context, accurate and timely short-run
flood forecasting—at lead times from one day to several days—
it plays an important role in early warning, evacuation planning, and
real-time reservoir and gate operation.
Traditional hydrological and hydraulic models rely on
Process-based representations of rainfall–runoff, river routing,
and dynamics in storage, but their implementation can be limited.
Constrained by data scarcity, calibration complexity, and sensitivity to
non-stationary climatic and land-use conditions. Data-driven
ML- and deep learning-based approaches
have therefore emerged as attractive alternatives or comple-
ments, learning non-linear relationships directly from histor-
-physical data without explicit physical parameterization. Within
This paradigm, Support Vector Machines (SVMs), and Long
The LSTM neural networks represent two
opposing modelling philosophies: margin-based kernel clas-
sifiers versus sequence-aware recurrent networks.
SVMs have been successfully used in various hydrological applications.
such applications as rainfall estimation and flood susceptibility.
It includes mapping and water quality assessment because of its toughness.
on relatively small, high-dimensional datasets. As, LSTMs
Specialized variety of recurrent neural network, have demonstrated
strated strong capabilities in modeling time-dependent pro-
cesses including streamflow forecasting, rainfall–runoff sim-
ulation, flood hydrograph prediction. Recent works that
incorporate attention mechanisms, vector-direction enhance-
ments, or physical constraints report high predictive skill for
time-series of flood-related, but often at the cost of increased
complexity and data requirements.
While much progress has been made, a pressing need remains for sys-
Tematic comparative studies that evaluate classical ML models.
pitting classical kernel methods, for instance, SVMs against deep sequence models, LSTMs
under the same experimental pipeline, using realistic
hydro-meteorological datasets. This paper addresses that need.
by doing a controlled comparison between SVM and LSTM for
short-term prediction of flood occurrence by using Flood Risks in
India dataset from Kaggle, which integrates meteorological data,
hydrological, and socio-environmental attributes. Core
question is how much additional benefit an LSTM provides.
over an SVM when both models operate on the same 7-
day windows of rainfall, river discharge, temperature, and
humidity.
The main contributions of this work are three-fold. First,
A reproducible pipeline is thus intended to transform raw daily
hydro-meteorological data into supervised sequences suitable
for both SVM and LSTM while maintaining temporal order through
Chronological train–test splitting. Second, a detailed empirical
The comparison is presented through accuracy, precision, recall.
F1-score, confusion matrices, and ROC–AUC, giving a
nuanced view of model behavior beyond a single aggregate
metric. Third, the operational implications of the observed
trade-off between high recall and high false-alarm rate are
Highlight, through analysis, the limitation of the current feature.
setting and outlining directions to improve flood forecasting.
Performance. 

\section{Related Work}
Early data-driven hydrological modeling relied on methods
This includes multilayer perceptrons and radial basis function networks.
Regression trees and early formulations of SVM, which demonstrate
strated that complex non-linear rainfall–runoff relationships
can be approximated without explicit physical models when
There are sufficient historical data available. These methods have
have been used to predict river discharge, water levels, and flooding.
occurrence, competitive performances compared to
from calibrated conceptual models in data-rich basins. However
their capabilities of capturing long-range temporal dependencies, and
rapidly changing hydrological regimes has been limited by the
Lack of explicit mechanisms for memory.
LSTM networks overcome some of these limitations by
introducing the gated memory cells that can selectively retain or
forget information over multiple time steps, making them well
suited for time-series forecasting problems. In flood forecast-
ing, LSTMs have been applied for the prediction of river discharge at
various lead times, to reconstruct the flood hydrographs, and to
anticipate peak flows using combinations of rainfall, previ-,
ous discharge and other hydro-meteorological inputs. Studies
have announced high performances measured by metrics such as
Nash–Sutcliffe efficiency (NSE), Kling–Gupta efficiency
(KGE), especially when LSTMs are augmented with attention
mechanisms or physically-inspired input representations.
Hybrid and Physics-informed deep learning models bring together
LSTMs combined with hydrological process models leveraging both
physical knowledge and data-driven flexibility in enhancing ro-
business and generalization. Such models can correct system-
approach conceptual model biases, downscale coarse predictions,
or fuse multi-source data including satellite rainfall and in-
Situ gauge measurements. While powerful, these approaches
-typically require significant model engineering and high-
quality data, which may limit their applicability in data-
scare catchments.
SVMs have maintained relevance due to their strong the-
theoretical foundations and practicality on moderate-sized tabu-
lar datasets. In flood-related studies, the SVMs have been used
for flood susceptibility mapping, and classification of flood-prone
zones, and rainfall or runoff prediction, often with quite good
performance for relatively simple feature sets, when combined
bined with appropriate kernels and feature selection tech-
niques. Comparative studies have sometimes found that while
deep networks outperform SVMs for complex spatio-temporal
prediction tasks, SVMs may be competitive or even superior.
when data are limited or when features are already informative.
and engineered.
A recent systematic review of time-series data mining
Flood prediction underlines the important evaluation of
models using several complementary metrics, such as accuracy,
F1-score, RMSE, NSE, ROC–AUC) and of the careful handling
class imbalance in event prediction tasks. This work
Contributes to this body of work by comparing an RBF-kernel
SVM and a small LSTM architecture sharing the same pipeline.
on the Flood Risk in India dataset, focusing on
on a short-term binary flood occurrence prediction, rather than
continuous discharge forecasting. 

\section{Dataset and Study Area}
These experiments use the dataset Flood Risk in India.
are available on Kaggle and represent synthetic yet realistic
records made to simulate the different hydro-meteorological and
socio-environmental conditions across regions in India. The full
dataset contains 10,000 records with 14 columns, including
latitude, longitude, rainfall (mm), temperature (◦C), humidity
(%), river discharge (m3/s), water level (m), elevation (m), land
cover, soil type, population density, infrastructure, historical
flood indicator, and a binary variable indicating whether a
FLOOD OCCURRED. These characteristics altogether capture climatic
forcing, catchment response, topographic controls and anthro-
Pogenic exposure relevant to flood processes.
By applying the modeling pipeline, this study identifies
and maps key columns to standardized variable names: rain-
fall mm, max temp C, humidity %, river discharge m3s,
and flood label. A synthetic date column is generated
, from 1 January 2019 onwards so that each row corre-
sponds to one daily observation, yielding a processed dataset
with six fields: Date, Rainfall mm, River Discharge m3s,
Max Temp C, Humidity %, and Flood Label. Descriptive
Statistics show that the rainfall ranges approximately from
0.014 mm to 299.97 mm, discharge from about 0.04 m3/ s to
4,999.70 m3/s, temperature from about 15.00 ◦C to 44.99 ◦C,
and humidity from about 20.00% to 99.99%, which reflects a wide
variety of hydrological regimes.
The flood label is strictly binary, with the processed dataset
consisting of 5,057 flood samples and 4,943 non-flood samples.
This corresponds to 50.57% of days when flood happens.
Balanced distribution differs from many real-world time se-
ries, where flood days are relatively rare, but simplifies the
Comparative evaluation of classification algorithms without
extreme imbalance effects. It is important to note that, while
the spatial context is abstracted in the synthetic dataset.
the hydro-meteorological ranges are broadly consistent with
Indian climatic and riverine conditions provide a dataset with
useful benchmark for methodological studies.
\begin{table}[htbp]
\caption{Summary Statistics of Key Variables (10{,}000 Records)}
\centering
\begin{tabular}{lrrrr}
\toprule
\textbf{Variable} & \textbf{Min} & \textbf{Mean} & \textbf{Max} & \textbf{Std} \\
\midrule
Rainfall\_mm & 0.01 & 150.02 & 299.97 & 86.03 \\
River\_Disch.\,(m\(^3\)/s) & 0.04 & 2515.72 & 4999.70 & 1441.71 \\
Max\_Temp\_C & 15.00 & 29.96 & 44.99 & 8.67 \\
Humidity\_\% & 20.00 & 59.75 & 99.99 & 23.14 \\
Flood\_Label & 0.00 & 0.51 & 1.00 & 0.50 \\
\bottomrule
\end{tabular}
\label{tab:stats}
\end{table}

% ==========================================================
% ===============  FIGURE PLACEHOLDER: DATAVIZ  ============
% ==========  REPLACE WITH FINAL FIGURE  =======
% ==========================================================


\begin{figure}[h!]
    \centering
    \includegraphics[width=0.9\linewidth]{data_vis_1.png}
    \label{fig:datavis1}
    \includegraphics[width=0.9\linewidth]{data_vis_2.png}
    \caption{Data Visualization}
    \label{fig:datavis2}
\end{figure}


\section{Methodology}
\subsection{Problem Formulation}
The objective is to perform short-term binary classification of daily flood occurrence using recent hydro-meteorological observations. For each day \(t\), the model receives a sequence of input vectors \(\mathbf{x}_\tau \in \mathbb{R}^4\) representing rainfall, river discharge, maximum temperature, and humidity over the previous seven days, and predicts whether a flood occurs on the next day. Formally, the input sequence is
\[
\mathbf{X}_t = [\mathbf{x}_{t-6},\mathbf{x}_{t-5},\dots,\mathbf{x}_t]
\]
with \(\mathbf{X}_t \in \mathbb{R}^{7 \times 4}\), and the target label \(y_{t+1} \in \{0,1\}\) denotes flood or no flood on day \(t+1\). The modeling task is to learn a mapping
\[
f : \mathbb{R}^{7 \times 4} \rightarrow \{0,1\}
\]
such that \(f(\mathbf{X}_t)\) approximates \(y_{t+1}\) for unseen sequences. 

\subsection{Data Preprocessing and Sequence Construction}
The raw CSV file is first loaded, and relevant columns
are automatically detected based on string patterns matching
rainfall, temperature, humidity, river discharge, and flood oc-
currence, which are then renamed to standardized forms. The
Flood label is processed to ensure strict binarization and mapping
any non-zero entries to 1 and verifying the resulting unique
value set {0, 1}. A synthetic Date column is created by using a
daily frequency starting from 2019-01-01, so that each record
corresponds to a unique day in the constructed time series.
From this preprocessed dataset, a target dataset of supervised learning is
built using a sliding window with lookback L= 7 days.
For each index i from 6 to N − 1, where N = 10,000 is the
total number of records, the input sequence comprises the four
predictors from days i − 6 through i, and the corresponding
target is the flood label at day i. This procedure yields a 3D
Input array of shape (9,993, 7, 4) and a label vector of length
9,993, since the first six samples cannot form a complete 7-day
window.
To simulate a real forecasting situation and avoid
temporal leakage, the sequences are chronologically split into
80% training and 20% testing subsets. More specifically, the first
7,994 sequences are used for training and the remaining 1,999
sequences are reserved for evaluation, with both sets maintain-
ing approximately balanced flood/no-flood distributions. This
Temporal split reflects the realistic constraint that the future
The outcomes should not affect model training. 

\begin{table}[htbp]
\caption{Train--Test Split and Class Distribution for Sequences.}
\centering
\begin{tabular}{lrrr}
\toprule
\textbf{Set} & \textbf{Samples} & \textbf{Flood (1)} & \textbf{No Flood (0)} \\
\midrule
Training & 7{,}994 & 4{,}053 (50.70\%) & 3{,}941 (49.30\%) \\
Test     & 1{,}999 & 1{,}000 (50.03\%) & 999 (49.97\%) \\
\bottomrule
\end{tabular}
\label{tab:split}
\end{table}

Feature scaling is applied using \texttt{MinMaxScaler} from scikit-learn fitted on the training data only, transforming each feature \(x\) according to
\[
x_{\text{scaled}} = \frac{x - x_{\min}^{\text{train}}}{x_{\max}^{\text{train}} - x_{\min}^{\text{train}}},
\]
where \(x_{\min}^{\text{train}}\) and \(x_{\max}^{\text{train}}\) are the minimum and maximum observed values in the training set. The same scaler is then applied to the test set, ensuring consistent feature ranges without leaking test information into training. 

\subsection{Handling Class Imbalance}
Although the dataset is nearly balanced, the pipeline explicitly computes class weights to make the methodology robust to potential imbalance and to reflect the higher importance of correctly identifying flood days. The class weights are computed using \texttt{compute\_class\_weight} from scikit-learn, yielding weights approximately 1.0142 for class 0 (no flood) and 0.9862 for class 1 (flood) based on the training distribution. These weights are passed as \texttt{class\_weight='balanced'} to the SVM and as per-sample weights to the LSTM, thereby slightly emphasizing minority-class errors in the loss function. 

\begin{table}[htbp]
\caption{Computed Class Weights from Training Data.}
\centering
\begin{tabular}{lrr}
\toprule
\textbf{Class} & \textbf{Count} & \textbf{Weight} \\
\midrule
0 (No Flood) & 3{,}941 & 1.0142 \\
1 (Flood)    & 4{,}053 & 0.9862 \\
\bottomrule
\end{tabular}
\label{tab:weights}
\end{table}

\subsection{Code Methodology and Implementation Workflow}
To ensure reproducibility and clarity, the entire modeling process is implemented in a structured, modular Python notebook that follows a logical flow closely resembling a block-diagram style workflow. The key stages of the code methodology are summarized below, corresponding to the major print banners in the console output. 

\begin{enumerate}
    \item \textbf{Environment Setup and Imports:} The code installs and imports \texttt{tensorflow}, \texttt{scikit-learn}, \texttt{pandas}, \texttt{numpy}, \texttt{matplotlib}, and \texttt{seaborn}, and prints their versions to confirm the runtime environment.
    \item \textbf{Dataset Selection and Loading:} The user selects between uploading a custom CSV or generating a synthetic dataset; in this study, the uploaded file \texttt{flood\_risk\_dataset\_india.csv} is chosen, and its shape and column names are printed.
    \item \textbf{Column Mapping and Basic Processing:} The code automatically maps dataset-specific column names to standardized internal names (e.g., Rainfall\_mm, River\_Discharge\_m3s, Max\_Temp\_C, Humidity\_\%, Flood\_Label), converts the flood label to binary, and constructs the Date column and a compact processed DataFrame.
    \item \textbf{Exploratory Data Visualization:} Time-series plots, histograms, and distribution plots of the key variables and flood labels are generated using \texttt{matplotlib} and \texttt{seaborn}, providing visual insight into ranges, variability, and class balance.
    \item \textbf{Sequence Construction:} A dedicated function constructs 7-day sliding windows for the four predictors and aligns the corresponding flood label as the target, producing the 3D input tensor and 1D label vector used by both models.
    \item \textbf{Chronological Train--Test Split:} The sequences are split chronologically into training and test sets, and the composition of each set is printed, including the number and percentage of flood and no-flood samples.
    \item \textbf{Feature Normalization:} A \texttt{MinMaxScaler} is fit on the training data and applied to both training and test sets, with the transformed ranges and prevention of data leakage explicitly documented in the console output.
    \item \textbf{Class Weight and Sample Weight Computation:} Class frequencies in the training set are used to compute class weights, and a vector of sample weights is generated to be passed to the LSTM training routine.
    \item \textbf{SVM Model Training:} The 3D input arrays are flattened to 2D (7 days \(\times\) 4 features), an RBF-kernel \texttt{SVC} model is instantiated with probabilistic outputs and balanced class weights, and the model is trained and evaluated on the test set.
    \item \textbf{LSTM Model Definition and Training:} A Keras Sequential model is constructed with an LSTM layer, dropout, dense layers, and a sigmoid output, then compiled with binary cross-entropy loss and Adam optimizer, and trained for up to 50 epochs with early stopping and sample weights.
    \item \textbf{Training History Visualization:} The training and validation loss and accuracy curves are plotted, revealing convergence behavior and potential overfitting or underfitting patterns for the LSTM.
    \item \textbf{Model Evaluation and Comparison:} Both models are evaluated via classification reports, confusion matrices, ROC curves, and a combined metrics table, which are then visualized using heatmaps, ROC plots, and bar charts.
    \item \textbf{Model Serialization:} The trained SVM, LSTM, and the fitted scaler are saved to disk as \texttt{svm\_flood\_model.pkl}, \texttt{lstm\_flood\_model.h5}, and \texttt{scaler.pkl}, enabling reuse and deployment.
\end{enumerate}

% ==========================================================
% ===============  FIGURE PLACEHOLDER: CODEFLOW  ============
% ==========  REPLACE WITH FINAL FIGURE  =======
% ==========================================================


\begin{figure}[htbp]
    \centering
    \includegraphics[width=0.9\linewidth]{flow.png}
    \caption{Flow diagram summarizing the code workflow from data loading, preprocessing, and sequence construction through SVM/LSTM training, evaluation, and model saving.}
    \label{fig:codeflow}
\end{figure}


\section{Model Architectures}
\subsection{Support Vector Machine}
The SVM classifier is implemented using scikit-learn's \texttt{SVC} with a radial basis function (RBF) kernel, which enables modeling non-linear decision boundaries between flood and non-flood classes. The 3D input sequences are flattened into 28-dimensional vectors (7 days \(\times\) 4 features) so that each sample can be processed as a fixed-length feature vector. The SVM is configured with \(C=1.0\), \(\gamma=\texttt{scale}\), \texttt{probability=True}, and \texttt{class\_weight='balanced'}, allowing both probabilistic predictions and automatic scaling of the kernel width based on the data. 

This configuration exploits the SVM's ability to maximize the margin between classes in feature space while handling moderate feature dimensionality without explicit feature engineering. However, by flattening the sequences, temporal order is encoded only implicitly, and the model cannot natively capture sequential dependencies beyond the fixed positional encoding determined by feature index. 

\subsection{Long Short-Term Memory Network}
The LSTM model is constructed in TensorFlow/Keras as a sequential network tailored to exploit the full 7-day temporal structure of the inputs. The architecture consists of an LSTM layer with 64 units that ingests sequences of shape \((7,4)\), followed by a dropout layer with rate 0.2, a dense layer of 32 units with ReLU activation, another dropout layer, and a final dense layer with a single neuron and sigmoid activation for binary classification. The total number of trainable parameters is 19{,}777, as reported by the model summary. 

The model is compiled using the Adam optimizer with binary cross-entropy loss and accuracy as a monitored metric, and is trained for up to 50 epochs with a batch size of 32 and 20\% of the training data held out for validation. An early stopping callback monitors validation loss and restores the best model if the loss does not improve for 10 consecutive epochs, which in this experiment results in training stopping after 12 epochs and restoring weights from epoch 2. Sample weights based on class frequencies are passed during training to slightly emphasize flood samples in the loss function. 

% ==========================================================
% ===============  FIGURE PLACEHOLDER: LSTM_SUMMARY  ============
% ==========  REPLACE WITH FINAL FIGURE  =======
% ==========================================================


\begin{figure}[htbp]
    \centering

    % First image (top)
    \includegraphics[width=0.85\linewidth]{L1.png}

    % Second image (bottom)
    \includegraphics[width=0.85\linewidth]{l2.png}

    \caption{LSTM model summary showing layers and parameter counts.}
    \label{fig:lstm_summary}
\end{figure}


\section{Experimental Setup}
Both models are trained and evaluated in a Python environment with TensorFlow~2.19.0, pandas~2.2.2, and NumPy~2.0.2, as reported at the beginning of the console output. The same chronological train--test split, normalized features, and binary labels are used for SVM and LSTM to ensure a fair comparison. All evaluations are performed on the fixed 1{,}999-sample test set, and random seeds are kept at their default values, implying that small variations could occur across runs but the qualitative behavior remains similar. 

Performance is assessed using standard classification metrics suitable for binary event prediction in hydrology: accuracy, precision, recall, F1-score, confusion matrix, and ROC--AUC. For the positive class (flood), precision is defined as \(\text{Precision} = \frac{TP}{TP+FP}\), recall as \(\text{Recall} = \frac{TP}{TP+FN}\), accuracy as \(\text{Accuracy} = \frac{TP+TN}{TP+TN+FP+FN}\), and the F1-score as the harmonic mean of precision and recall, \(F1 = 2\cdot\frac{\text{Precision} \cdot \text{Recall}}{\text{Precision}+\text{Recall}}\). ROC--AUC quantifies the probability that a randomly chosen flood day receives a higher predicted probability than a randomly chosen non-flood day, averaged over all decision thresholds. 

\section{Results}
\subsection{LSTM Training Behavior}
The LSTM training history indicates that training and validation accuracies fluctuate around 0.49--0.51, while training and validation losses stay close to 0.693 throughout the 12 epochs before early stopping. This behavior is characteristic of a classifier whose output probabilities remain near 0.5 for most samples, suggesting difficulty in learning strongly discriminative decision boundaries under the chosen architecture and feature set. The early stopping mechanism restores the model weights from epoch 2, implying that later epochs did not provide consistent improvements in validation loss and potentially preventing overfitting to noise. 

% ==========================================================
% ===============  FIGURE PLACEHOLDER: LSTM_HISTORY  ============
% ==========  REPLACE WITH FINAL FIGURE  =======
% ==========================================================


\begin{figure}[htbp]
    \centering
    
    % First image (Loss plot)
    \includegraphics[width=0.9\linewidth]{eva.png}
    \vspace{0.7em} % space between plots

    % Second image (Accuracy plot)
    \includegraphics[width=0.9\linewidth]{eva2.png}

    \caption{LSTM training history showing training and validation loss (top) and accuracy (bottom) across epochs.}
    \label{fig:lstm_history}
\end{figure}



\subsection{Quantitative Performance Comparison}
On the 1{,}999-sample test set, the SVM model produces 964 predicted flood samples and 1{,}035 predicted non-flood samples. The resulting classification report shows that for the no-flood class (0), precision is 0.4986, recall is 0.5165, and F1-score is 0.5074, while for the flood class (1), precision is 0.4990, recall is 0.4810, and F1-score is 0.4898. Overall accuracy is 0.4987, macro-averaged F1-score is 0.4986, and ROC--AUC is 0.5069, indicating near-random performance but with slightly better-than-chance discrimination. 

The LSTM model, by contrast, predicts 1{,}964 flood samples out of 1{,}999 test instances, effectively labeling almost all days as floods. Its classification report shows that for the no-flood class (0), precision is 0.6000, recall is 0.0210, and F1-score is 0.0406, while for the flood class (1), precision is 0.5020, recall is 0.9860, and F1-score is 0.6653. Overall accuracy is 0.5038, macro-averaged F1-score is 0.3530, and ROC--AUC is 0.5013, indicating very strong sensitivity for flood events but poor specificity and essentially random threshold-independent discrimination. 

\begin{table}[htbp]
\caption{SVM and LSTM Test-Set Performance Metrics.}
\centering
\begin{tabular}{lrrrrr}
\toprule
\textbf{Model} & \textbf{Accuracy} & \textbf{Precision} & \textbf{Recall} & \textbf{F1} & \textbf{ROC--AUC} \\
\midrule
SVM  & 0.4987 & 0.4990 & 0.4810 & 0.4898 & 0.5069 \\
LSTM & 0.5038 & 0.5020 & 0.9860 & 0.6653 & 0.5013 \\
\bottomrule
\end{tabular}
\label{tab:metrics}
\end{table}

\subsection{Confusion Matrices}
The SVM confusion matrix on the test set is
\[
\begin{bmatrix}
516 & 483 \\
519 & 481
\end{bmatrix},
\]
where the first row corresponds to true no-flood days and the second row to true flood days. This matrix indicates that the SVM correctly identifies 516 out of 999 no-flood days and 481 out of 1{,}000 flood days, while misclassifying 483 no-flood days as floods and 519 flood days as no-flood. The fairly symmetric structure reflects the model's near-random behavior, with a slight tendency to correctly classify non-flood days. 

The LSTM confusion matrix is
\[
\begin{bmatrix}
21 & 978 \\
14 & 986
\end{bmatrix},
\]
showing that the model correctly identifies 986 out of 1{,}000 flood days but only 21 out of 999 no-flood days. This extreme skew toward predicting floods results in a very small number of missed floods (14 false negatives) but a very large number of false alarms (978 false positives), which explains the high recall but low precision and poor ROC--AUC. 

% ==========================================================
% ===============  FIGURE PLACEHOLDER: CONFUSION  ============
% ==========  REPLACE WITH FINAL FIGURE  =======
% ==========================================================


\begin{figure}[htbp]
    \centering
    \includegraphics[width=0.9\linewidth]{download.png}
    \caption{Heatmap of the model confusion matrix on the test set.}
    \label{fig:confusion}
\end{figure}


\subsection{ROC Curves and Bar-Chart Comparison}
ROC curves constructed for both models show that their points lie only slightly above the diagonal line corresponding to a random classifier. The SVM ROC curve achieves an AUC of 0.5069, while the LSTM ROC curve yields an AUC of 0.5013, confirming that neither model extracts strong threshold-independent discrimination from the available 7-day feature set. The slight advantage of the SVM in ROC--AUC is overshadowed by its much lower F1-score compared to the LSTM. 

A bar chart comparing accuracy, precision, recall, F1-score, and ROC--AUC across the two models illustrates that the LSTM outperforms the SVM in four out of five metrics (accuracy, precision, recall, F1-score), while the SVM is slightly superior only in ROC--AUC. Given the focus on flood-event detection, the LSTM is identified as the overall winner based on F1-score, with a reported improvement of 35.83\% relative to the SVM's F1-score. 

% ==========================================================
% ===============  FIGURE PLACEHOLDER: ROC  ============
% ==========  REPLACE WITH FINAL FIGURE  =======
% ==========================================================


\begin{figure}[htbp]
    \centering
    \includegraphics[width=0.9\linewidth]{roc.png}
    \caption{ROC curve comparison between SVM and LSTM models with corresponding AUC values.}
    \label{fig:roc}
\end{figure}


% ==========================================================
% ===============  FIGURE PLACEHOLDER: BAR  ============
% ==========  REPLACE WITH FINAL FIGURE  =======
% ==========================================================


\begin{figure}[htbp]
    \centering
    \includegraphics[width=0.9\linewidth]{last.png}
    \caption{Performance comparison of SVM and LSTM models across accuracy, precision, recall, F1-score, and ROC--AUC metrics.}
    \label{fig:bar}
\end{figure}


\section{Discussion}
The comparative results reveal a clear trade-off between
sensitivity of flood events and discriminative capability in general
as quantified by ROC–AUC. The SVM behaves like a near-
random classifier, offering only marginal separation between
flood and no-flood days regardless of using a non-linear ker-
lsm with balanced class weights. The LSTM, in contrast,
Aggressively predicts the flood class, with very few false
negatives, but at the cost of a very high false-positive rate and
poor ROC–AUC.
Operationally speaking, LSTM is doing the following:
Could be acceptable in situations where missing a flood is
much costlier than issuing a false alarm, such as in critical
Infrastructure protection or high-risk communities. However,
persistent over-warning will result in warning fatigue, reduced
trust in the system, and unnecessary economic costs, which
suggests that calibration of the decision threshold and explicit
control of false-alarm rates is essential for deployment. The
nearly random ROC–AUC values of both models would suggest that
under the current set of features and architecture, the underlying
data are inadequate to produce a usefully sharp signal.
discrimination across a range of thresholds.
Several factors can be involved in the limited performance.
observed. First, only four dynamic predictors are used, and key
hydrological states such as antecedent soil moisture, catch-
storage of ment, and detailed basin topology are not explicitly
represented, which may be crucial to distinguish accurately between
: (1) A total of 59 features are used in the classification scheme for disinguishing flood and non-flood conditions. (2) The LSTM
its architecture is relatively shallow and not extensively tuned;
deeper or attention-enhanced architectures may capture
more subtle temporal patterns if appropriate regularization is
applied. Third, synthetic nature of the dataset may impose
limits to how well models can generalize physical behaviours.
In particular, if the generative assumptions are far from real
catchment dynamics.
For the SVM, performance might be improved by con-
structing engineered features from the 7-day window, such as
cumulative rainfall, maximum discharge or slope of discharge,
which could provide more hydrologically interpretable signals
as opposed to raw daily values. Hyperparameter tuning of C and γ,
exploring alternative kernels, incorporating static attributes.
(e.g., elevation, soil type) could also enhance its discriminative
power. In the case of the LSTM, tests can be run with different lookback
windows, dropout rates, learning rates and loss functions e.g.,
Whereas, focal loss may help strike a better balance between recall.
, including precision, without collapsing into near-constant flooding.
Prediction.
 

\section{Conclusion and Future Work}
This paper presented a comparative analysis of SVM and LSTM models for short-term flood occurrence prediction from hydro-meteorological time series derived from the \emph{Flood Risk in India} dataset. A unified preprocessing pipeline was designed to map raw CSV data into a temporally ordered dataset, construct 7-day sliding input windows, perform chronological train--test splitting, normalize features, compute class weights, and train both an RBF-kernel SVM and a 64-unit LSTM network. On a 1{,}999-sample test set, using accuracy, precision, recall, F1-score, confusion matrices, and ROC--AUC, the LSTM was found to substantially outperform the SVM on F1-scores and recall for flood events, with both showing near-random ROC--AUC values, indicating poor threshold-independent discriminability. 

These results indicate that the LSTM strongly favors sensitivity to floods, eliminating almost all false negatives but at the cost of a very high number of false positives, while the SVM provides more balanced but weak predictions. These trade-offs highlight the need for richer hydro-meteorological and catchment descriptors, more advanced sequence modeling architectures, and explicit threshold calibration when designing flood early-warning systems based on data-driven models. In future work, the integration of additional variables is necessary, such as water level, soil moisture proxies, or static basin attributes. Furthermore, this work will also apply attention-based and physics-informed LSTM variants and apply the pipeline to real gauge-based datasets from specific catchments in order to evaluate physically meaningful performance metrics such as NSE and KGE.

\section*{Acknowledgment} The authors want to thank the contributors who made the data \emph{Flood Risk in India} available at Kaggle and the open-source communities that developed TensorFlow, scikit-learn, pandas, NumPy, Matplotlib, and Seaborn, whose software was employed in this study.


\begin{thebibliography}{00}
\bibitem{ref1}
A. Mosavi, P. Ozturk and K.-W. Chau, “Flood prediction using machine learning models: A review,” \textit{Water}, vol. 10, no. 11, p. 1536, 2018.

\bibitem{ref2}
F. Kratzert, D. Klotz, M. Herrnegger, A. K. Sampson, S. Hochreiter and G. S. Nearing, 
“Toward improved predictions in ungauged basins: Combining LSTM networks and hydrologic models,” 
\textit{Water Resources Research}, vol. 55, no. 12, pp. 11377–11402, 2019.

\bibitem{ref3}
Y. Zhang and R. S. Govindaraju, 
“Prediction of watershed runoff using Bayesian neural networks,” 
\textit{Water Resources Research}, vol. 36, no. 3, pp. 753–762, 2000.

\bibitem{ref4}
M. K. Tiwari and C. Chatterjee, 
“Development of an accurate and reliable hourly flood forecasting model using wavelet-bootstrap-ANN hybrid approach,” 
\textit{Journal of Hydrology}, vol. 394, pp. 458–470, 2010.

\bibitem{ref5}
J. Zhang, Y. Zhu, X. Zhang, M. Ye and J. Yang, 
“Short-term flood forecasting using LSTM networks with attention mechanism,” 
\textit{Journal of Hydrology}, vol. 597, 126277, 2021.

\bibitem{ref6}
M. Sit, B. Z. Demiray, Z. Xiang, G. J. Ewing, Y. Sermet and I. Demir, 
“A comprehensive review of LSTM networks for hydrologic modeling,” 
\textit{Engineering Applications of Artificial Intelligence}, vol. 93, 103852, 2020.

\bibitem{ref7}
Kaggle, “Flood Risk in India Dataset,” 2023. [Online]. Available: https://www.kaggle.com  
\end{thebibliography}

\end{document}